\section{In the Beginning}

\begin{quotex}
In the beginning God \flright{\textit{Genesis} 1:1}

\end{quotex}

God is whatever was at the beginning. For the physicist, that is some ill-defined quantity of energy confined to a
vanishingly small space. It is in some sense transcendent because it does not include any of the fundamental particles
that define matter; that happens later as the universe begins to cool.

Yet, that inchoate pre-matter contains all the physical laws and fundamental constants that allow the formation of
stars, galaxies, planet, life, human beings, and so on. Hence, the physicist is not an atheist, he just does not
recognize God.

\paragraph{Revelation}
Physics can explain atoms, molecules, even crystals, but not a living thing. Living things have structures apart from
their physical and chemical components. No single physical property accounts for life.

But what is structure but its form? So when a metaphysician claims that the soul is the form or the body, that is what
is meant. Therefore the form, or the soul, is likewise real.

Human life is inexplicable by physics. Consciousness is not hidden in any theory, and there is nothing to suggest that
intelligent beings, capable of understanding the laws of physics, should arise.

The human being discovers that there is moral universe with its own laws, unknown and unknowable to physics. Moreover,
unlike stars and planets, the human being wonders why he exists and what his purpose is. This can only be known by
revelation.

Moses, who is more important than any scientist, passed on what was revealed to him.

\begin{itemize}
\item God is a person, not an inchoate cloud of mass/energy. 
\item God has a plan for us to be with Him. 
\item God has provided us with guidance to create a good society and a path to lead us to Him. 
\end{itemize}

Obviously, that is simplified, but still captures the essential. Unfortunately, the principles of guidance are called
“commandments”, a word which has taken on misleading connotations over the centuries. Commandments are often
misunderstood as directives to compel us to do what we don't want to do, with the result that we will be
punished for failure to follow such commandments.

Rather, commandments like forbidding murder, theft, adultery, covetousness, seem obvious to create both a peaceable
social structure and an harmonious inner life. Who would be surprised that breaking those commandments leads to social
conflict and inner turmoil? So we are not arbitrarily “punished” for breaking commandments, rather we are warned of the
adverse consequences of doing so.

It is analogous to warning your young child not to touch the hot stove, lest she hurt herself. The “hurt” is not your
punishment, but rather the inevitable consequence.

We learn to love God's commandments for that reason; if we fail, God loves and forgives. Think in those
terms and you won't be left feeling dark and somber about your future prospects.

\flright{\itshape Posted on 2022-08-14 by Cologero}